\section{Limitações}
\label{sec:limitacoes}

O escopo impõe limitações necessárias. O protótipo de bancada não reproduz
atmosfera explosiva, bomba certificada, vazão, interferências de pista nem todos
os tempos de uma operação comercial. Atuadores comuns representam visualmente
o bloqueio, mas não demonstram aptidão industrial. A consulta real depende de
aprovação externa, contratos e documentação técnica ainda não disponíveis; por
isso, o simulador é a fonte padrão do desenvolvimento.

A posição dos bicos exibida na planta React é simulada no navegador. O ACK do
ESP32 informa processamento do comando e estado lógico da saída GPIO, não
deslocamento do solenoide, posição do pino ou presença do bico no suporte. Como
a bancada atual não possui microswitch nem sensor Hall, desenergizar uma saída
com o bico já retirado não o devolve nem comprova novo bloqueio mecânico.

O MVP de software também é restrito ao próprio computador. Identidades de
frentista e supervisor são autodeclaradas, e o token de supervisor é global;
portanto, os testes atuais não demonstram autenticação individual, segregação de
funções ou segurança de uma implantação em rede. O uso de HMAC reduz a exposição
de identificadores na persistência e no cache, mas não transforma o dado de
origem em anônimo nem elimina a necessidade de proteger o segredo.

O estudo de caso pode conter poucos incidentes, registros heterogêneos e custos
incompletos. Entrevistas intencionais fornecem compreensão contextual, mas não
estimativa populacional. A presença do pesquisador na organização pode influir
na coleta e interpretação. A análise econômica será contrafactual e sensível a
volume de consultas, tarifa da GCC, vida útil e percentual de prevenção. Mesmo
que H1 seja atendida em bancada, isso não comprovará eficácia, segurança ou
conformidade em uma instalação real.

O projeto limita-se a uma trava independente no suporte do bico. Não foram
modelados nem avaliados comando de bombeamento, controle de vazão, integração
com controlador de pista ou alterações na automação da empresa.

Por fim, a descrição de combustível no cadastro veicular pode não representar
alterações posteriores, conversões ou particularidades não cobertas pela
política. Casos ambíguos serão bloqueados para revisão, o que preserva segurança
mas pode reduzir disponibilidade e aumentar o tempo de atendimento.
